\chapter{Introduction}
\label{chap:introduction}
Riding a bicycle is a common activity for humans, but designing a controller to autonomously stabilize a bicycle presents significant engineering challenges.
Traditional model-based controllers, such as \gls{lqr} and \gls{mpc}, require an accurate system model.
However, some parameters of a bicycle, i.e., the center of gravity, can be difficult to determine and may fluctuate due to varying loads or additional passengers.
In this project, we aim to design and implement a \gls{ddmpc} to stabilize an autonomous mini-bike.
The method does not identify the system model of the bike and directly uses the collected data to design the controller. The goal is to stabilize the bike at a constant speed and the bike does not fall down.
Besides, we aim to keep both the lean angle and steering angle within defined safe limits for safety.


To balance a bicycle, a torque should be generated to prevent falling on its side due to gravity. To achieve this, different techniques for torque generation can be introduced such as adding a device to control the leaning, or by controlling the steering angle of the bicycle while moving. While the bicycle is stationary, a flywheel can generate the required torque to achieve the balancing \cite{kawaguchi2011}. Another approach is by using a control moment gyroscope (CMG) which was tested out by He and Zhao \cite{he2015}. Since controlling the steering angle to produce the balancing torque does not need additional devices, it was considered by numerous researchers and tested against emerging control algorithms. Balancing using the steering is based on steering into the direction of falling. A drawback of using only the steering angle is that the bicycle needs an initial speed before being able to balance.

Before implementing a control strategy for balancing a bicycle, a rigorous system model must be established for most techniques, though some approaches are model-free. System models can be classified as either kinematic or dynamic. The kinematic model does not consider the applied forces on the system, while the dynamic model considers the torques and forces interacting with the system. Researchers set different parameters for the dynamic models depending on the bicycle under study. Some parameters could be neglected to simplify the system. Complex systems emerging from accurate modelling demand sophisticated control algorithms, which require more computation power. Examples of dynamic models for the bicycle include the Whipple model \cite{whipple1899stability} and the point mass model \cite{limebeer2006}. Andersson et al. \cite{andersson2019} proposed a robust proportional-integral-derivative (PID) to control a bicycle based on the point mass model.

Since researchers neglect some parameters in the bicycle models to simplify computations, some uncertainties arise. Therefore, a robust controller can mitigate the errors generated by such models. Robust controllers include the \gls{smc}, which was proposed by Chen and Dao \cite{chen2010}, and the \gls{lqr}, which was proposed by Anjumol and Jisha \cite{anjumol2014}. The controllers proposed in the literature are largely tested in simulations rather than on a real bicycle.




\section{Structure of thesis}
This thesis is structured as follows:
Chapter 1 provides an introduction, including the background, motivation, and objectives of the research.
Chapter 2 discusses the theoretical foundations, focusing on \gls{mpc} and related concepts.
Chapter 3 describes the simulation setup and the implementation of the \gls{ddmpc} approach.
Chapter 4 presents the experimental results, including the setup and data collection from the mini-bike.
\ref{chap}Chapter 5 concludes the report and discusses future work.

